% !TEX root = somerville-sample.tex

\input somerville-font-setup.tex

\setupbodyfontenvironment[default][em=italic]

\setupSomerville
  {Width=97,Weight=400,Optical Size=12,Math Weight=700,Arrow Length=200,Arrow Head=100,Integral Slant=0}
  {somerville-body}
% Separate variable-font instance for heavier text and math in headings.

\setupSomerville
  {Width=100,Weight=700,Optical Size=24,Math Weight=900,Arrow Length=200,Arrow Head=100,Integral Slant=0}
  {somerville-headings}

\setuphead[title]
  [style={\switchtobodyfont[somerville-headings]\tfc}]
\setuphead[subject]
  [style={\switchtobodyfont[somerville-headings]\tfb}]

\setupbodyfont[somerville-body,12pt]

\starttext
\title{Somerville Sample: \m{E=mc^2}}
\blank[small]

Here is some text {\it some italics} and {\bf some bold} and {\bi some bold italics}

Comparisons

{\it text italics} \m{math italics} \crlf
{\bf text bold} \m{\bf math bold} \crlf
{\bi text bolditalic} \m{\bi math bolditalic} \crlf


{\it ABCDEF} \m{ABCDEF} \crlf
{\bf ABCDEF} \m{\bf ABCDEF} \crlf
{\bi ABCDEF} \m{\bi ABCDEF} \crlf

\m{x<\mathbb R + \mathbb T}

\m{\mathrm{A} {\mathbb A} {\mathrm B} {\mathbb B} \mathrm{R} {\mathbb R} R}


For \m{a = b = l} \m{f'u'=g'=h'=u'=y'} \m{a,b,c\in\mathbb{R}}, the identity \m{a^2+b^2=c^2} sits
alongside fractions \m{\frac{1}{2}} and radicals \m{\sqrt{2}}.
Nested scripts: \m{|x| = x_i^2+x_{i_j}^{n+1}+e^{x^2}}.

\subject{Fractions and radicals: \m{\sqrt{1+x^2}}}
\startformula
  x=\frac{-b\pm\sqrt{b^2-4ac}}{2a}
  \qquad
  \frac{1}{1+\frac{1}{1+x}}
  \qquad
  \sqrt{1+\sqrt{1+\sqrt{x}}}
\stopformula

\subject{Integrals, sums, and products: \m{\int{f(x)\,\mathrm{d}x}_{a}^{b}}}
\startformula
  \int{x^2\,\mathrm{d}x}_{0}^{1}=\frac{1}{3}
  \qquad
  \int{e^{-x^2}\,\mathrm{d}x}_{-\infty}^{\infty}=\sqrt{p}
\stopformula
\startformula
  \sum{k}_{k=1}^{n}=\frac{n(n+1)}{2}
  \qquad
  \prod{\frac{k+1}{k}}_{k=1}^{n}=n+1
\stopformula

\subject{Delimiters and matrices: \m{\det\mathbf{A}}}
\startformula
  \left(\frac{x+1}{x-1}\right)^2
  \qquad
  \left\langle\frac{u}{\|u\|},\frac{v}{\|v\|}\right\rangle
  \qquad
  A=\left(\startmatrix
    \NC a \NC b \NR
    \NC c \NC d \NR
  \stopmatrix\right)
\stopformula

\subject{Limits, accents, and arrows: \m{\lim_{h\to0}\frac{f(x+h)-f(x)}{h}}}
\startformula
  \lim_{h\to0}\frac{f(x+h)-f(x)}{h}=f'(x)
  \qquad
  \widehat{f}(x)=\int{f(x)e^{-2p i xx}\,\mathrm{d}x}_{-\infty}^{\infty}
\stopformula
\startformula
  u'+v'=w'
  \qquad
  \overline{z}=x-iy
  \qquad
  x\in A\cap B\iff x\in A\ \land\ x\in B
\stopformula


\subject{Integral sizes and italic correction: \m{\int{}_{a}^{b}}}
Each row requests a larger integral at the same body size.
Left: side limits. Right: limits above and below.
The identical limit letters make horizontal shifts easier to compare.
% ConTeXt's size sets both the requested height and depth.
% The font's available variants determine the actual glyph selected.
% Change INSL in the axis settings to compare upright and slanted integrals.

\unexpanded\def\IntegralCorrectionTest#1{%
  \startformula
    \text{size #1}\qquad
    \int[size=#1,method=nolimits]{f(x) \mathrm{d}x}_{a}^{a}
    \qquad
    \int[size=#1,method=limits]{f(x) \mathrm{d}x}_{a}^{a}
  \stopformula
}
\m{\int{f(x) dx}_{a}^{a}}
\IntegralCorrectionTest{6pt}
\IntegralCorrectionTest{9pt}
\IntegralCorrectionTest{12pt}
\IntegralCorrectionTest{18pt}
\IntegralCorrectionTest{36pt}
\IntegralCorrectionTest{48pt}

\page


\subject{Miscellanea \emph{mathematica}: \m{e^{ip}+1=0}}

Lorem ipsum dolor sit amet, \m{x\in\mathbb{R}} consectetur adipiscing elit.
Sed vestibulum \m{[0,1]}, ut \m{e>0} fermentum \emph{aculis ornare}.
Curabitur consequat \m{a_n=1/(n+1)} neque vitae sollicitudin;
pellentesque habitant \m{b_n=(-1)^n a_n} morbi tristique senectus et netus.
Praesent \m{N\geq1}, commodo viverra \m{s_N=\sum{a_n}_{n=1}^{N}} quisque feugiat
lobortis sapien. Nulla facilisi: etiam \m{d>0} tincidunt augue,
vel \m{M<\infty} porttitor massa, donec malesuada hendrerit arcu.
Suspendisse finibus ipsum nec turpis sodales, non posuere justo dictum.
\startformula
  \int{x^n\,\mathrm{d}x}_{0}^{1}=\frac{1}{n+1},
  \qquad
  \sum{\frac{z^n}{n!}}_{n=0}^{\infty}=e^z,
  \qquad
  \lim_{n\to\infty}\left(1+\frac{t}{n}\right)^n=e^t.
\stopformula

Vestibulum ante ipsum primis in faucibus, \m{f} mollis nunc
et \m{f'(a)\neq0} luctus sapien. Donec posuere \m{y=f(a)+f'(a)(x-a)} ultricies
ligula, vel \emph{fermentum arcu dignissim sed}. Maecenas
\m{R_2(h)=O(h^3)} pretium lectus apud \m{h=0}. Inter
\m{a<b} et \m{g\leq d}, sollicitudin ornare nibh,
quisque consequat \m{\frac{p}{q}} vitae facilisis
\m{\sqrt{1+\sqrt{1+x^2}}}. Sed posuere, nunc et semper, consequat
urna consectetur, aliquam erat volutpat, etiam porta tellus.
\startformula
  f(a+h)=f(a)+hf'(a)+\frac{h^2}{2!}f''(a)
  +\frac{1}{2}\int{(h-t)^2f'''(a+t)\,\mathrm{d}t}_{0}^{h}.
\stopformula

Integer pretium \m{(m,n)\in\mathbb{Z}^2} vestibulum magna, dictum
viverra \m{w_{m,n}=e^{-(m^2+n^2)}} semper aliquam.
Praesent commodo nunc vel nulla malesuada, et consectetur
purus elementum. Hic \m{|z|<1}, illic
\m{\mathop{\mathrm{Re}}(s)>1}, atque interdum sollicitudin
\m{\prod{(1+u_k)}_{k=1}^{N}} mauris vitae risus accumsan.
Phasellus \m{A\subseteq B\subseteq X} porttitor ornare;
\m{A\cup B}, \m{A\cap B}, et \m{X\setminus A} habitant platea.
\startformula
  \sum_{m=-\infty}^{\infty} \sum{e^{-(m^2+n^2)}}_{n=-\infty}^{\infty}
  =\left(\sum{e^{-k^2}}_{k=-\infty}^{\infty}\right)^2,
  \qquad
  \sum{z^k}_{k=0}^{N}=\frac{1-z^{N+1}}{1-z}.
\stopformula

Aliquam erat volutpat. Sit \m{V=\mathbb{R}^3}, etiam
\m{u,v\in V}, atque \m{\langle u,v\rangle=\sum{u_jv_j}_{j=1}^{3}} nunc vehicula.
Nulla \m{\|u\|=\sqrt{\langle u,u\rangle}} facilisis eleifend,
sed \m{t} tempus fermentum. Vestibulum
\m{A=(a_{ij})_{i,j=1}^3} consequat nibh, ac malesuada lorem.
Cum \m{A^{\mathrm{T}}=A}, dignissim ipsum; cum
\m{A^{\mathrm{T}}A=I}, porttitor arcu. Mauris consequat
sapien vitae turpis cursus, vestibulum blandit massa dapibus.
\startformula
  \left(\startmatrix
    \NC 2 \NC -1 \NC 0 \NR
    \NC -1 \NC 2 \NC -1 \NR
    \NC 0 \NC -1 \NC 2 \NR
  \stopmatrix\right)
  \left(\startmatrix
    \NC x_1 \NR \NC x_2 \NR \NC x_3 \NR
  \stopmatrix\right)
  =\left(\startmatrix
    \NC 2x_1-x_2 \NR
    \NC -x_1+2x_2-x_3 \NR
    \NC -x_2+2x_3 \NR
  \stopmatrix\right).
\stopformula

Suspendisse potenti, \m{r(x)=e^{-x^2/2}/\sqrt{2p}} faucibus
mollis ante vestibulum. Etiam \m{m=0}, semper
\m{s^2=1}, et \m{r>0} tincidunt quam
vitae posuere sapien. Nam \m{X},
\m{\mathbb{E}[X]} et \m{\mathbb{P}(X\leq r)} gravida
ligula praesent. Donec dignissim blandit massa,
ab \m{-\infty} ad \m{+\infty}, commodo arcu
pellentesque habitant morbi tristique senectus et netus.
\startformula
  \int{r(x)\,\mathrm{d}x}_{-\infty}^{\infty}=1,
  \qquad
  \int{xr(x)\,\mathrm{d}x}_{-\infty}^{\infty}=0,
  \qquad
  \int{x^2r(x)\,\mathrm{d}x}_{-\infty}^{\infty}=1.
\stopformula

Phasellus consectetur \m{\sin(nx)} et \m{\cos(nx)}
convallis ornare. Vestibulum \m{\widehat f(n)} malesuada
mauris, dum \m{e^{inx}=\cos(nx)+i\sin(nx)} vulputate
felis. Quisque \m{f(x)=\cos x+\frac12\sin(2x)}
et \m{2p}; etiam \m{g(x)=e^{-x^2}}
urna sodales, finibus nisl vitae, porta neque.
Ut enim ad minim veniam, cursus commodo facilisis consequat.
\startformula
  \widehat f(n)=\frac{1}{2p}\int{f(x)e^{-inx}\,\mathrm{d}x}_{-p}^{p},
  \qquad
  \frac{1}{2p}\int{|f(x)|^2\,\mathrm{d}x}_{-p}^{p}
    =\sum{|\widehat f(n)|^2}_{n\in\mathbb{Z}}.
\stopformula

Nunc varius, \m{u=u(x,t)} sollicitudin \m{\partial_tu=k\partial_x^2u}
cum \m{k>0}. Praesent tincidunt augue vel massa
blandit, malesuada arcu semper \m{t\to\infty}.
Etiam \m{\nabla f}, interdum \m{Df}, atque
\m{\partial^2f/(\partial x\,\partial y)} consequat
nibh vestibulum. Aenean \m{u_0\in L^1(\mathbb{R})}
finibus lorem, pellentesque sapien, dignissim semper magna.
\startformula
  u(x,t)=\frac{1}{\sqrt{4pk t}}
    \int{\exp\left(-\frac{(x-y)^2}{4k t}\right)
    u_0(y)\,\mathrm{d}y}_{-\infty}^{\infty},
  \qquad t>0.
\stopformula

Denique lorem ipsum, ornare minima: \m{x'}, \m{x''},
\m{x'''}, \m{a_{n+1}}, et \m{b^{k^2}_{i_j}}.
Mauris \m{\overline{z}}, vestibulum \m{\widehat{g}}, atque
porttitor \m{X\to Y} donec \m{Y\to Z}.
Sive \m{p\Rightarrow q} vel \m{p\Leftrightarrow q}, posuere
ligula \m{0<e\ll1}, malesuada cursus tellus,
etiam finibus \m{\sqrt{a^2+b^2}\leq|a|+|b|}.

\subject{De campis electromagneticis: \m{\nabla\times\mathbf{E}=-\partial_t\mathbf{B}}}

Lorem ipsum Maxwellianum, sit \m{\mathbf{E}=\mathbf{E}(\mathbf{x},t)}
et \m{\mathbf{B}=\mathbf{B}(\mathbf{x},t)} in spatio continuo.
Vestibulum \m{r}, praesent \m{\mathbf{J}}, atque
\m{c^{-2}=m_0e_0} consectetur adipiscing elit.
Curabitur fluxus, nunc et semper, per superficies orientatas procedit;
nulla facilisi, nam \m{\nabla\cdot\mathbf{B}=0} dignissim
magna et \m{\partial_tr+\nabla\cdot\mathbf{J}=0} consequat
sapien. Donec malesuada, felis vel porttitor gravida, arcu erat
sollicitudin lorem, vitae elementum risus neque sed justo.
\startformula
  \startalign[n=2]
    \NC \nabla\cdot\mathbf{E}=\frac{r}{e_0}
    \NC \nabla\cdot\mathbf{B}=0 \NR
    \NC \nabla\times\mathbf{E}=-\frac{\partial\mathbf{B}}{\partial t}
    \NC \nabla\times\mathbf{B}=m_0\mathbf{J}
      +m_0e_0\frac{\partial\mathbf{E}}{\partial t} \NR
  \stopalign
\stopformula

Praesent integralia \m{\partial V} et \m{\partial S} circumveniunt,
dum \m{\mathbf{n}} normalis, \m{\mathrm{d}\boldsymbol{\ell}} tangentialis,
et \m{Q_V=\int{r\,\mathrm{d}V}_{V}} massa verborum ornant.
Aenean quis lacus vitae sapien bibendum cursus. Integer ut augue
convallis, ornare arcu sed, dictum mauris; etiam \m{F_B}
interdum varius, sed \m{e_0>0} semper constans.
\startformula
  \int{\mathbf{E}\cdot\mathbf{n}\,\mathrm{d}S}_{\partial V}
    =\frac{1}{e_0}\int{r\,\mathrm{d}V}_{V},
  \qquad
  \oint{\mathbf{E}\cdot\mathrm{d}\boldsymbol{\ell}}_{\partial S}
    =-\frac{\mathrm{d}}{\mathrm{d}t}
      \int{\mathbf{B}\cdot\mathbf{n}\,\mathrm{d}S}_{S}.
\stopformula

Pellentesque Poyntingiana, \m{\mathbf{S}=m_0^{-1}\mathbf{E}\times\mathbf{B}}
vehicula lucis, et \m{u} condimentum energiae. Morbi tristique
senectus et netus, \m{\mathbf{F}=q(\mathbf{E}+\mathbf{v}\times\mathbf{B})}
fermentum ipsum. Sed ut perspiciatis unde omnis, vestibulum
\m{\mathbf{A}} atque \m{f} in faucibus orci luctus posuere.
\startformula
  u=\frac{e_0}{2}|\mathbf{E}|^2
    +\frac{1}{2m_0}|\mathbf{B}|^2,
  \qquad
  \frac{\partial u}{\partial t}+\nabla\cdot\mathbf{S}
    =-\mathbf{J}\cdot\mathbf{E}.
\stopformula
\startformula
  \mathbf{B}=\nabla\times\mathbf{A},
  \qquad
  \mathbf{E}=-\nabla f-\frac{\partial\mathbf{A}}{\partial t},
  \qquad
  \nabla^2\mathbf{E}-\frac{1}{c^2}
    \frac{\partial^2\mathbf{E}}{\partial t^2}=\mathbf{0}
  \quad(r=0,\ \mathbf{J}=\mathbf{0}).
\stopformula

\subject{De undis et quantis: \m{\widehat{H}q=Eq}}

Suspendisse Schrödingeriana, \m{q\in L^2(\mathbb{R}^3)}
volutpat magna, \m{\hbar>0} tincidunt minima.
Nam \m{\widehat{H}q=Eq} commodo lectus, sed
\m{\widehat{\mathbf{p}}=-i\hbar\nabla} dignissim arcu.
Vivamus ultricies, ipsum vitae malesuada tristique, urna erat
consequat nibh, sit amet vulputate erat nulla non sem.
Inter \m{\langle q,q\rangle=1} et \m{|q|^2}
vestibulum probabilitatis, quisque elementum pretium sapien.
\startformula
  i\hbar\frac{\partial q}{\partial t}
    =\left[-\frac{\hbar^2}{2m}\nabla^2+V(\mathbf{x},t)\right]q,
  \qquad
  \int{|q(\mathbf{x},t)|^2\,\mathrm{d}^3\mathbf{x}}_{\mathbb{R}^3}=1.
\stopformula
\startformula
  \left[\widehat{x},\widehat{p}\right]=i\hbar I,
  \qquad
  D x\,D p\geq\frac{\hbar}{2},
  \qquad
  E_n=\hbar w\left(n+\frac12\right),\quad n=0,1,2,\ldots.
\stopformula

Aliquam Pauli matrices, \m{\boldsymbol{s}=(s_1,s_2,s_3)}
ornare neque. Donec \m{s_j^2=I} tempus erat,
\m{s_1s_2=is_3} blandit justo, atque
\m{\mathbf{a}\cdot\boldsymbol{s}} porttitor risus.
Mauris at tellus et lectus facilisis, nec consequat nulla;
curabitur convallis nisl id metus tempor, vitae fringilla neque
sagittis. Nunc \m{\mathbf{a},\mathbf{b}\in\mathbb{R}^3} semper
interdum, et \m{\mathbf{a}\times\mathbf{b}} vestibulum ipsum.
\startformula
  s_1=\left(\startmatrix
    \NC 0 \NC 1 \NR \NC 1 \NC 0 \NR
  \stopmatrix\right),\quad
  s_2=\left(\startmatrix
    \NC 0 \NC -i \NR \NC i \NC 0 \NR
  \stopmatrix\right),\quad
  s_3=\left(\startmatrix
    \NC 1 \NC 0 \NR \NC 0 \NC -1 \NR
  \stopmatrix\right).
\stopformula
\startformula
  (\mathbf{a}\cdot\boldsymbol{s})
  (\mathbf{b}\cdot\boldsymbol{s})
    =(\mathbf{a}\cdot\mathbf{b})I
      +i(\mathbf{a}\times\mathbf{b})\cdot\boldsymbol{s}.
\stopformula

\subject{De geometria et gravitate: \m{R_{mn}-\frac12 Rg_{mn}}}

Aenean Einsteiniana, \m{g_{mn}} vestibulum curvatum,
et \m{R_{mn}} sollicitudin magna. Proin \m{L}
ornare constans, \m{T_{mn}} fermentum materiae, atque
\m{\nabla_m T^{mn}=0} consequat conservatum.
Lorem ipsum dolor sit amet, consectetur adipiscing elit;
morbi at arcu nec augue accumsan suscipit. Etiam
\m{\mathrm{d}s^2=g_{mn}\,\mathrm{d}x^m\mathrm{d}x^n}
interdum posuere, nec \m{t} neque \m{l} facilisis desunt.
\startformula
  R_{mn}-\frac12 Rg_{mn}+L g_{mn}
    =\frac{8p G}{c^4}T_{mn},
  \qquad
  E^2=|\mathbf{p}|^2c^2+m^2c^4.
\stopformula
\startformula
  \frac{\mathrm{d}^2x^m}{\mathrm{d}t^2}
    +G^m_{ab}
      \frac{\mathrm{d}x^a}{\mathrm{d}t}
      \frac{\mathrm{d}x^b}{\mathrm{d}t}=0,
  \qquad
  G^m_{ab}
    =\frac12 g^{mn}
      (\partial_a g_{nb}
       +\partial_b g_{na}
       -\partial_n g_{ab}).
\stopformula

\subject{De fluidis et variationibus: \m{\nabla\cdot\mathbf{u}=0}}

Vestibulum Navierianum, \m{\mathbf{u}=\mathbf{u}(\mathbf{x},t)}
fluvius continuus, \m{p} pretium pressionis, \m{n>0}
viscositas serena. Nullam tincidunt libero sed ipsum faucibus,
sed \m{\nabla\cdot\mathbf{u}=0} sagittis incompressibilis.
Etiam \m{r>0} constans et \m{\mathbf{f}} sollicitudin externa,
pellentesque habitant morbi tristique senectus. Donec at lectus
vitae justo posuere maximus, ut \m{t\geq0} ornare nunc et semper.
\startformula
  \frac{\partial\mathbf{u}}{\partial t}
    +(\mathbf{u}\cdot\nabla)\mathbf{u}
    =-\frac{1}{r}\nabla p+nD\mathbf{u}+\mathbf{f},
  \qquad
  \nabla\cdot\mathbf{u}=0.
\stopformula

Curabitur Lagrangiana, \m{L=L(\mathbf{q},\dot{\mathbf{q}},t)}
vestibulum motus, \m{S[\mathbf{q}]} condimentum actionis.
Sed \m{d\mathbf{q}(t_0)=d\mathbf{q}(t_1)=\mathbf{0}}
ornare fines, et \m{d S=0} praesent dictum.
Praesent non ligula nec ipsum consequat sodales; maecenas
\m{\mathbf{p}=\partial L/\partial\dot{\mathbf{q}}} fermentum
neque, dum \m{H=\mathbf{p}\cdot\dot{\mathbf{q}}-L} finibus arcu.
\startformula
  S[\mathbf{q}]=\int{L(\mathbf{q},\dot{\mathbf{q}},t)\,\mathrm{d}t}_{t_0}^{t_1},
  \qquad
  \frac{\mathrm{d}}{\mathrm{d}t}
    \frac{\partial L}{\partial\dot q_j}
    -\frac{\partial L}{\partial q_j}=0.
\stopformula
\startformula
  \dot q_j=\frac{\partial H}{\partial p_j},
  \qquad
  \dot p_j=-\frac{\partial H}{\partial q_j},
  \qquad
  \{f,g\}=\sum{
    \left(\frac{\partial f}{\partial q_j}\frac{\partial g}{\partial p_j}
    -\frac{\partial f}{\partial p_j}\frac{\partial g}{\partial q_j}\right)}.
\stopformula

\subject{De numeris et infinitis: \m{z(2)=\frac{p^2}{6}}}

Lorem Eulerianum, \m{e^{ip}+1=0} dictum celebre;
Riemanniana \m{z(s)} porta infinita, ubi
\m{\mathop{\mathrm{Re}}(s)>1} vestibulum absolutum.
Quisque \m{p\in\mathbb{P}} primus, \m{n\geq1} integer,
et \m{G(s)} mollis integralis. Fusce sit amet diam
at nisi consequat egestas, vitae congue lorem suscipit.
Donec \m{G(n+1)=n!} volutpat sapien, at
\m{z(2)=p^2/6} dignissim lectus.
\startformula
  z(s)=\sum{\frac{1}{n^s}}_{n=1}^{\infty}
    =\prod{\frac{1}{1-p^{-s}}}_{p\ \mathrm{prime}},
  \qquad
  G(s)=\int{t^{s-1}e^{-t}\,\mathrm{d}t}_{0}^{\infty}.
\stopformula
\startformula
  G(s)G(1-s)=\frac{p}{\sin p s},
  \qquad
  n!\sim\sqrt{2p n}\left(\frac{n}{e}\right)^n.
\stopformula

Morbi Cauchyana, \m{f} holomorpha, \m{g} curva positiva,
et \m{z_0} punctum interius. Vestibulum ante ipsum primis,
\m{f^{(n)}(z_0)} sollicitudin alta, \m{n!} consequat frequens.
Sed ut perspiciatis, finibus neque sed, blandit massa;
\m{\oint{f(z)\,\mathrm{d}z}_{g}} circumdat margines,
nec \m{2p i} nec \m{z-z_0} memoria deficit.
\startformula
  f^{(n)}(z_0)=\frac{n!}{2p i}
    \oint{\frac{f(z)}{(z-z_0)^{n+1}}\,\mathrm{d}z}_{g},
  \qquad
  \int{\frac{\mathrm{d}x}{1+x^2}}_{-\infty}^{\infty}=p.
\stopformula

\subject{De probabilitate et informatione: \m{\mathbb{E}[\mathbf{X}]=\boldsymbol{m}}}

Suspendisse Bayesiana, \m{\mathbb{P}(B)>0} porta aperta,
\m{\mathbb{P}(A\mid B)} condimentum novum. Praesent
\m{p_i>0}, \m{q_i>0}, et \m{\sum{p_i}_{i}=\sum{q_i}_{i}=1}
ornare simplex. Cras consectetur sapien vel urna commodo,
interdum \m{\boldsymbol{t}\in\mathbb{R}^d} fermentum variabile,
interdum \m{\nabla_{\boldsymbol{t}}\ell} sagittis ascendens.
Mauris accumsan lorem ut magna pretium, in placerat massa dictum.
\startformula
  \mathbb{P}(A\mid B)
    =\frac{\mathbb{P}(B\mid A)\mathbb{P}(A)}{\mathbb{P}(B)},
  \qquad
  H(p)=-\sum{p_i\log p_i}_{i},
  \qquad
  D(p\|q)=\sum{p_i\log\frac{p_i}{q_i}}_{i}\geq0.
\stopformula

Denique Gaussiana, \m{\mathbf{x},\boldsymbol{m}\in\mathbb{R}^d},
et \m{\boldsymbol{S}} symmetrica positiva. Pellentesque
\m{\det\boldsymbol{S}>0} dignissim plena, dum
\m{(\mathbf{x}-\boldsymbol{m})^{\mathrm{T}}} vestibulum
transpositum. Lorem ipsum dolor sit amet, consectetur adipiscing
elit; etiam dictum magna, porta serena, finibus purus.
\startformula
  p(\mathbf{x})=
    \frac{\exp\!\left[-\frac12
      (\mathbf{x}-\boldsymbol{m})^{\mathrm{T}}
      \boldsymbol{S}^{-1}
      (\mathbf{x}-\boldsymbol{m})\right]}
    {(2p)^{d/2}\sqrt{\det\boldsymbol{S}}}.
\stopformula

\subject{De fractionibus continuis et radicibus: \m{\sqrt[n]{1+x^n}}}

Lorem ipsum profundum, \m{\frac{1}{1+\frac{1}{1+x}}} vestibulum
angustum, \m{\sqrt{\frac{a_i^2+b_i^2}{1+c_i^2}}} porttitor altum.
Praesent \m{x^{\frac{1}{2}}}, \m{x_{\frac{1}{2}}}, atque
\m{e^{-\frac{x^2}{2s^2}}} sollicitudin minuta.
Curabitur superius et inferius, cras posuere nibh; nam
\m{\frac{\sqrt{1+x^2}}{\sqrt{1-x^2}}} inter verba procedit.
\startformula
  v=1+\frac{1}{1+\frac{1}{1+\frac{1}{1+\cdots}}},
  \qquad
  \sqrt{1+\sqrt{1+\sqrt{1+\sqrt{1+x}}}},
  \qquad
  \sqrt[\scriptstyle n+1]{\frac{1+x^{2n+2}}{1+x^2}}.
\stopformula
\startformula
  \frac{\displaystyle\sum{\frac{a_j}{\sqrt{1+b_j^2}}}_{j=1}^{n}}
    {\displaystyle 1+\int{\frac{\sqrt{1+t^4}}{1+t^2}\,\mathrm{d}t}_{0}^{1}}
  \qquad
  \sqrt{\frac{\displaystyle 1+\frac{x^2}{1+y^2}}
    {\displaystyle 1+\frac{y^2}{1+x^2}}}
  \qquad
  \sqrt[3]{\sqrt[5]{\sqrt[7]{1+z^{105}}}}.
\stopformula

\subject{De indicibus et litteris inclinatis: \m{\mathbf{T}_{i_j}^{k^\ell}}}

Vestibulum \m{f_i}, \m{f^j}, \m{f_i^j}, \m{f^{j^2}_{i_k}}
et \m{f_{i_{j_k}}^{n^{m^2}}} ornare creberrimum.
Aliquam \m{J_a^b}, \m{T_a^b}, \m{V_a^b}, \m{W_a^b},
\m{G_a^b} et \m{\int{}_{a}^{b}} finibus obliquis;
\m{f'^2}, \m{f''_i}, \m{f'''(x)}, \m{x^{\prime\prime}_n}
semper praesent. Donec \m{\mathbf{A}_{i_j}^{k^\ell}}
nec \m{\boldsymbol{S}_{a_b}^{g^d}}
facilisis desunt, in hac habitasse platea dictumst.
\startformula
  {}^{14}_{6}\mathrm{C}
  \qquad {}^{238}_{92}\mathrm{U}
  \qquad {}^{a}_{b}\mathbf{T}_{g}^{d}
  \qquad
  R^{r}{}_{smn}
  =\partial_mG^{r}_{ns}
   -\partial_nG^{r}_{ms}
   +G^{r}_{ml}G^{l}_{ns}
   -G^{r}_{nl}G^{l}_{ms}.
\stopformula
\startformula
  \sum{a_i^{\sum{b_{ij}}_{j=1}^{m}}}_{i=1}^{n}
  \qquad
  e^{\displaystyle-\frac{x^2}{2s^2}}
  \qquad
  \left.\frac{\mathrm{d}^{n}}{\mathrm{d}x^{n}}
    e^{-x^2}\right|_{x=0}
  \qquad
  \left(\frac{x_i^{j_k}}{y_{m_n}^{p^q}}\right)^{\!2}.
\stopformula

\subject{De limitibus multiplicibus: \m{\sum{a_{ij}}_{i,j=1}^{n}}}

Integer \m{\sum{}_{i=1}^{n}}, \m{\prod{}_{j=1}^{m}}, atque
\m{\int{}_{0}^{1}} vehicula parva; sed in spatio amplo
\m{\bigcup_{a\in A}U_a} et
\m{\bigcap_{n\geq1}F_n} porttitor magna.
Praesent \m{\lim_{n\to\infty}} volutpat lenta,
\m{\sup_{x\in K}|f(x)|} dignissim finita.

Etiam \m{\inf_{y\in Y}d(x,y)} sollicitudin propinqua.
\startformula
  \sum{a_{ij}}_{\startsubstack 1\leq i,j\leq n \NR i\neq j \stopsubstack}
  \qquad
  \prod{\left(1-\frac{1}{p^s}\right)^{-1}}_{\startsubstack p\leq N \NR p\in\mathbb{P} \stopsubstack}
  \qquad
  \lim_{\startsubstack n\to\infty \NR h\to0^+ \stopsubstack}
    \frac{f_n(x+h)-f_n(x)}{h}.
\stopformula
\startformula
  \bigcup_{n=1}^{\infty}\bigcap_{m=n}^{\infty}A_m
    \subseteq
  \bigcap_{n=1}^{\infty}\bigcup_{m=n}^{\infty}A_m,
  \qquad
  \left\|\sum{\mathbf{v}_j}_{j=1}^{n}\right\|^2
    =\sum{\langle\mathbf{v}_i,\mathbf{v}_j\rangle}_{i,j=1}^{n}.
\stopformula

\subject{De integralibus altissimis: \m{\int{\mathrm{d}u}_{M}}}

Aenean Stokesiana, \m{w} forma lenis, \m{M} varietas orientata,
et \m{\partial M} terminus ordinatus. Nulla facilisi:
\m{\mathrm{d}^2w=0}, \m{\nabla\times\nabla f=\mathbf{0}},
et \m{\nabla\cdot(\nabla\times\mathbf{A})=0} semper concordant.
Sed \m{\int{\mathrm{d}w}_{M}=\int{w}_{\partial M}}
vestibulum antiquum; praesent limites superiores atque inferiores,
commodo arcu, finibus serenis.
\startformula
  \int[method=nolimits]{\frac{1}{\sqrt{1+\frac{x^2}{1+x^4}}}\,\mathrm{d}x}_{0}^{1}
  \qquad
  \int[method=limits]{\frac{1}{\sqrt{1+\frac{x^2}{1+x^4}}}\,\mathrm{d}x}_{0}^{1}.
\stopformula
\startformula
  \int[method=nolimits]{\left[\frac{\displaystyle\int{e^{-x^2t^2}\,\mathrm{d}t}_{0}^{1}}
      {\displaystyle 1+\frac{x^2}{1+x^2}}\right]
      e^{-x^2}\,\mathrm{d}x}_{-\infty}^{+\infty}
  \qquad
  \int[size=36pt,method=limits]{\frac{e^{-x^2}}{\sqrt{1+x^2}}\,\mathrm{d}x}_{-\infty}^{+\infty}.
\stopformula
\startformula
  \iint{(\nabla\times\mathbf{F})\cdot\mathbf{n}\,\mathrm{d}S}_{S}
    =\oint{\mathbf{F}\cdot\mathrm{d}\boldsymbol{\ell}}_{\partial S},
  \qquad
  \iiint{\nabla\cdot\mathbf{F}\,\mathrm{d}V}_{V}
    =\iint{\mathbf{F}\cdot\mathbf{n}\,\mathrm{d}S}_{\partial V}.
\stopformula

\subject{De accentibus et vinculis: \m{\widehat{\mathbf{A}\mathbf{B}}}}

Mauris \m{\hat f}, \m{\widehat{fg}}, \m{\widehat{fgh}},
\m{\widetilde{xyz}}, et \m{\overline{z_1+z_2}} condimentum
superpositum. Vestibulum \m{\dot x}, \m{\ddot x},
\m{\dot{\mathbf{x}}}, \m{\ddot{\mathbf{x}}}, atque
\m{\hat{\hat f}} ornare duplex. Phasellus volutpat
\m{\widehat{f_i^j}} tincidunt ardua, nec
\m{\overline{\overline{z}}} finibus obscuris.
\startformula
  \widehat{x}
  \quad \widehat{xy}
  \quad \widehat{xyz}
  \quad \widehat{x_1+x_2+x_3+x_4}
  \qquad
  \widetilde{\mathbf{A}\mathbf{B}\mathbf{C}}
  \quad
  \overline{\frac{z_1+z_2}{1+z_1z_2}}.
\stopformula
\startformula
  \overbrace{x_1+x_2+\cdots+x_n}^{n}
    =\underbrace{x+\cdots+x}_{n}\quad(x_j=x),
  \qquad
  \underbrace{\overbrace{a+b}^{u}+\overbrace{c+d}^{v}}_{u+v}.
\stopformula
\startformula
  \overrightarrow{AB}+\overrightarrow{BC}=\overrightarrow{AC},
  \qquad
  \overline{\underbrace{z_1+\cdots+z_n}_{n}}
    =\overline{z_1}+\cdots+\overline{z_n}.
\stopformula

\subject{De parenthesibus et matricibus magnis: \m{\mathbf{A}^{-1}\mathbf{A}=\mathbf{I}}}

Curabitur \m{\langle\mathbf{x},\mathbf{y}\rangle} spatium internum,
\m{\lfloor x\rfloor} gradus inferior, \m{\lceil x\rceil} superior.
Etiam \m{\{x\in\mathbb{R}:x^2<2\}} habitant aperte,
\m{|\det\mathbf{A}|} dignissim absoluta. In hac habitasse
platea dictumst, cras \m{\mathbf{A}^{-1}} sollicitudin inversa,
sed \m{\det\mathbf{A}\neq0} vestibulum necessarium.
\startformula
  \left\langle
    \left(\frac{1}{1+x^2},\frac{x}{1+x^2}\right),
    \left(\sqrt{\frac{1+y}{1-y}},\sqrt{\frac{1-y}{1+y}}\right)
  \right\rangle
  \qquad
  \left\lceil\frac{\displaystyle\sum{k^2}_{k=1}^{n}}
    {\displaystyle\sum{k}_{k=1}^{n}}\right\rceil.
\stopformula
\startformula
  \mathbf{H}_4=
  \left(\startmatrix
    \NC 1 \NC \frac12 \NC \frac13 \NC \frac14 \NR
    \NC \frac12 \NC \frac13 \NC \frac14 \NC \frac15 \NR
    \NC \frac13 \NC \frac14 \NC \frac15 \NC \frac16 \NR
    \NC \frac14 \NC \frac15 \NC \frac16 \NC \frac17 \NR
  \stopmatrix\right),
  \qquad
  \det\left(\startmatrix
    \NC 1 \NC x \NC x^2 \NR
    \NC 1 \NC y \NC y^2 \NR
    \NC 1 \NC z \NC z^2 \NR
  \stopmatrix\right)
  =(y-x)(z-x)(z-y).
\stopformula
\startformula
  \left\{\left[
    \left(\frac{\displaystyle\int{t^{a-1}(1-t)^{b-1}
      \,\mathrm{d}t}_{0}^{1}}{\displaystyle\frac{G(a)G(b)}
      {G(a+b)}}\right)
    \right]^2\right\}^{\!1/2}=1,
  \qquad a,b>0.
\stopformula

\subject{De casibus et aequationibus ordinatis: \m{d_{ij}}}

Lorem ipsum distributionum, \m{d_{ij}} porttitor discretum,
\m{d(x)} sollicitudin singularis. Praesent
\m{\int{d(x-a)f(x)\,\mathrm{d}x}_{\mathbb{R}}=f(a)}
volutpat concentrata, dum \m{H'(x)=d(x)} sensu distributionum.
Vivamus at \m{x=0} posuere media, etiam \m{x<0} et \m{x>0}
facilisis opposita. Donec ultricies sapien vitae massa cursus,
interdum condimentum tellus, semper vestibulum purus.
\startformula
  H(x)=\startcases
    \NC 0 \NC x<0, \NR
    \NC \frac12 \NC x=0, \NR
    \NC 1 \NC x>0, \NR
  \stopcases
  \qquad
  d_{ij}=\startcases
    \NC 1 \NC i=j, \NR
    \NC 0 \NC i\neq j. \NR
  \stopcases
\stopformula
\startformula
  \startalign[n=2,align={right,left}]
    \NC \|\mathbf{x}+\mathbf{y}\|^2
      \NC =\langle\mathbf{x}+\mathbf{y},\mathbf{x}+\mathbf{y}\rangle \NR
    \NC \NC =\|\mathbf{x}\|^2+2\langle\mathbf{x},\mathbf{y}\rangle
          +\|\mathbf{y}\|^2 \NR
    \NC \NC \leq\|\mathbf{x}\|^2+2\|\mathbf{x}\|\,\|\mathbf{y}\|
          +\|\mathbf{y}\|^2 \NR
    \NC \NC =(\|\mathbf{x}\|+\|\mathbf{y}\|)^2. \NR
  \stopalign
\stopformula

\subject{De formis exterioribus et identitatibus: \m{\mathrm{d}^2w=0}}

Aenean \m{w=\sum{a_i\,\mathrm{d}x^i}_{i}} vestibulum covarians,
\m{h=\sum{b_j\,\mathrm{d}x^j}_{j}} dignissim alternans.

Morbi \m{w\wedge h=-h\wedge w} commodo
antisymmetricum, et \m{\mathrm{d}^2=0} dictum breve.
Nullam \m{\mathcal{F}}, \m{\mathcal{L}}, \m{\mathcal{H}},
et \m{\mathbb{C}\subset\mathbb{H}} ornare varia;
\m{\mathbf{a}\otimes\mathbf{b}} condimentum tensoriale.
\startformula
\sum{\bigoplus{\frac{\bigotimes{a_n}}{\bigodot{b_n}}}_{i=1}^{n}} \neq
\bigoplus{\prod{\frac{\sum{a_n}}{\bigoplus{b_n}}}_{i=1}^{n}} 
\stopformula

Denique \m{\displaystyle\frac{\int{f(x)\,\mathrm{d}x}_{0}^{1}}{\sqrt{1+a^2}}}
inter verba surgit, dum \m{\scriptstyle\frac{a_i^2}{1+b_j^2}}
minuitur, et \m{\scriptscriptstyle x_{i_j}^{k^\ell}} paulatim recedit.
Lorem ipsum dolor sit amet, consectetur adipiscing elit; \m{\mathbf{0}},
\m{\mathbf{I}}, \m{\boldsymbol{a}}, \m{\boldsymbol{b}},
et \m{\boldsymbol{\nabla}} vestibulum novissimum.
Quisque \m{\left\|\frac{\mathbf{x}}{\|\mathbf{x}\|}\right\|=1}
finibus clausis, sollicitudin integra, commodo sereno.



\subject{De partitionibus et spectris: \m{Z_N(b)}}

Vestibulum harmonicum, \m{b>0}, \m{w_j>0}, et
\m{E_{\mathbf{n}}=\sum{\hbar w_j(n_j+\frac12)}_{j=1}^{N}}
fermentum discretum. Praesent \m{\mathbf{n}\in\mathbb{N}_0^N}
porttitor occupationis; quisque modus suam frequentiam retinet,
dum \m{Z_N(b)} consequat multiplicativa colligit.
\startformula
  Z_N(b)
  =\sum[size=24pt,method=limits]{\exp\!\left[
      -b\sum{\hbar w_j\left(n_j+\frac12\right)}_{j=1}^{N}
    \right]}_{\mathbf{n}\in\mathbb{N}_0^N}
  =\prod[size=24pt,method=limits]{\frac{e^{-b\hbar w_j/2}}
      {1-e^{-b\hbar w_j}}}_{j=1}^{N}.
\stopformula

Mauris \m{\mathbf{A}\in\mathbb{C}^{N\times N}} vestibulum spectrale,
\m{l_1,\ldots,l_N} radices ordinatae.
Nam \m{|z|\,\|\mathbf{A}\|<1} porta convergentiae,
et \m{\mathop{\mathrm{tr}}(\mathbf{A}^r)=\sum{l_j^r}_{j}}
dictum invarians. Integer finibus purus, curabitur series
concordes, donec determinantia in exponentia conveniunt.
\startformula
  \det(\mathbf{I}-z\mathbf{A})^{-1}
  =\prod[size=18pt,method=limits]{\frac{1}{1-zl_j}}_{j=1}^{N}
  =\exp\!\left[
    \sum[size=24pt,method=limits]{\frac{z^r}{r}\mathop{\mathrm{tr}}(\mathbf{A}^r)}_{r=1}^{\infty}
  \right].
\stopformula

\stoptext
